%=========================================================================
%  CHE F311 - Kinetics & Reactor Design
%  PROBLEMS AND SOLUTIONS
%
%  Every numerical question or example posed in the nine uploaded decks
%  (Slides/), quoted from the deck that sets it and solved with the method
%  and data the deck teaches.  Parts:
%    A  Ch1-krd (1).pdf   mole balances and rates            Q1-Q2
%    B  Ch2-krd (1).pdf   conversion and reactor sizing      Q3-Q8
%    C  Ch3-krd.pdf       rate laws and stoichiometry        Q9-Q14
%    D  Ch4-_Krd.pdf      isothermal design                  Q15-Q21
%    E  Ch5_-_krd.pdf     collection and analysis of rate data Q22-Q35
%    F  Ch6_krd.pdf       multiple reactions                 Q36-Q37
%    G  Ch7-_krd.pdf      non-isothermal design              Q38-Q42
%    H  Catalysis_Krd.pdf catalysis and pellets              Q43-Q44
%    I  Non_ideal_KRD.pdf residence time distributions       Q45-Q46
%
%  Every derived number is reproduced by krd-answers-check.py in this
%  folder.  Where a slide leaves a datum undefined the assumption is
%  stated in a Notes box instead of being filled in silently.
%=========================================================================
\documentclass[10pt,a4paper]{article}

\usepackage[a4paper,left=19mm,right=19mm,top=15mm,bottom=17mm]{geometry}
\usepackage[T1]{fontenc}
\usepackage{amsmath,amssymb,mathtools}
\usepackage{xcolor}
\usepackage{booktabs,tabularx,array}
\usepackage{enumitem}
\usepackage{fancyhdr}
\usepackage{tcolorbox}
\tcbuselibrary{breakable,skins}
\usepackage{siunitx}
\usepackage[colorlinks=true,linkcolor=accent,urlcolor=accent2]{hyperref}

\definecolor{accent}{HTML}{3B2A6B}
\definecolor{accent2}{HTML}{6B5BA6}
\definecolor{qframe}{HTML}{1B5E20}
\definecolor{qbg}{HTML}{EAF5EC}
\definecolor{mframe}{HTML}{1D6FA5}
\definecolor{mbg}{HTML}{EAF1F8}
\definecolor{boxbg}{HTML}{F4F1F7}
\definecolor{rule}{HTML}{A99FC4}
\definecolor{inbox}{HTML}{1A1430}
\definecolor{notecol}{HTML}{4E5260}
\definecolor{warncol}{HTML}{8E2B0F}
\definecolor{warnbg}{HTML}{FDF2EC}

\sisetup{detect-all,per-mode=symbol}
\emergencystretch=2em

\newtcolorbox{qbox}{breakable,enhanced,colback=qbg,colframe=qframe!65,
  boxrule=0.6pt,arc=1.5pt,left=6pt,right=6pt,top=4pt,bottom=4pt,
  before skip=5pt,after skip=4pt,fontupper=\small}
\newtcolorbox{mthbox}{breakable,enhanced,colback=mbg,colframe=mframe!65,
  boxrule=0.6pt,arc=1.5pt,left=6pt,right=6pt,top=4pt,bottom=4pt,
  before skip=5pt,after skip=4pt,fontupper=\small}
\newtcolorbox{ansbox}{enhanced,colback=boxbg,colframe=accent,boxrule=0.9pt,
  arc=1.5pt,left=6pt,right=6pt,top=5pt,bottom=5pt,before skip=5pt,after skip=6pt,
  fontupper=\small\color{inbox}}
\newtcolorbox{warnbox}{breakable,enhanced,colback=warnbg,colframe=warncol!60,
  boxrule=0.6pt,arc=1.5pt,left=6pt,right=6pt,top=4pt,bottom=4pt,
  before skip=5pt,after skip=4pt,fontupper=\small}

\newcommand{\spsection}[1]{%
  \par\addvspace{10pt}%
  \noindent\begin{tcolorbox}[colback=accent,colframe=accent,boxrule=0pt,arc=1.2pt,
      left=7pt,right=7pt,top=5pt,bottom=5pt,width=\linewidth,before skip=0pt,after skip=7pt]
    \color{white}\bfseries\large #1
  \end{tcolorbox}\par\vspace{-2pt}\nopagebreak}
\newcommand{\qtitle}[2]{%
  \par\addvspace{9pt}\noindent
  \begin{tcolorbox}[colback=accent2!12,colframe=accent2,boxrule=0.8pt,arc=1.5pt,
      left=7pt,right=7pt,top=4pt,bottom=4pt,before skip=0pt,after skip=4pt,width=\linewidth]
    {\large\bfseries\color{accent}#1}\hfill{\small\color{accent2}#2}
  \end{tcolorbox}\par\vspace{1pt}}
\newcommand{\src}[1]{\par\vspace{1pt}\noindent{\small\color{notecol}\textbf{Source:} #1}\par\vspace{3pt}}
\newcommand{\note}[1]{\par\vspace{1pt}\noindent{\small\color{notecol}#1}\par\vspace{2pt}}

\pagestyle{fancy}
\fancyhf{}
\fancyhead[L]{\small CHE F311 · Problems \& Solutions}
\fancyhead[R]{\small\thepage}
\renewcommand{\headrulewidth}{0.4pt}

\begin{document}

\noindent{\Large\bfseries\color{accent} KINETICS \& REACTOR DESIGN}\\
{\normalsize\color{accent2} CHE F311 · Problems and Solutions · I-Semester 2026-27}\\[4pt]
{\small Every numerical question in the uploaded decks, with its solution.}\\[6pt]

\noindent{\small Every question below is quoted from one of the uploaded files, and every
solution uses the method taught in the deck that set it.  The questions are grouped by
source deck (Parts A--I).  Where a slide prints its own answer the solution reproduces it;
where a slide leaves something undefined, the Notes box says what was assumed.  All numbers
are re-computed by \texttt{krd-answers-check.py} kept beside this file.}

\spsection{Question inventory}
\begin{footnotesize}\setlength{\tabcolsep}{4pt}
\begin{tabularx}{\linewidth}{@{}l l X r@{}}
\toprule
ID & Source & Topic & Solved at\\
\midrule
Q1  & Ch1, slide 7  & Rocket engine: rates of reaction from mole balance & p.~\pageref{q:1}\\
Q2  & Ch1, slide 12 & $A\to B$: CSTR, PFR and batch volumes/times, three rate laws & p.~\pageref{q:2}\\
Q3  & Ch2, slides 19--24 & Levenspiel sizing: CSTR and PFR at 80\% & p.~\pageref{q:3}\\
Q4  & Ch2, slides 31--36 & Two CSTRs in series (40\%, 80\%) & p.~\pageref{q:4}\\
Q5  & Ch2, slides 37--38 & Two PFRs in series & p.~\pageref{q:5}\\
Q6  & Ch2, slides 41--42 & Mixed reactor, unknown stoichiometry & p.~\pageref{q:6}\\
Q7  & Ch2, slides 45--46 & PFRs in parallel branches: feed split & p.~\pageref{q:7}\\
Q8  & Ch2, slide 59 & P2-7 adiabatic $A+B\to C$ data, (a)--(e) & p.~\pageref{q:8}\\
Q9  & Ch3, slides 5--7 & Relative rates: $2\mathrm{NO}+\mathrm{O_2}\to2\mathrm{NO_2}$ & p.~\pageref{q:9}\\
Q10 & Ch3, slides 6--7 & Relative rates: $2A+3B\to5C$ & p.~\pageref{q:10}\\
Q11 & Ch3, slides 42--44 & Saponification stoichiometric table & p.~\pageref{q:11}\\
Q12 & Ch3, slides 48--51 & Flow-system stoichiometric table, $A+2B\to C+D$ & p.~\pageref{q:12}\\
Q13 & Ch3 s.61 (+Ch5 s.77) & Dimerisation $2A\to R$ in MFR: rate equation & p.~\pageref{q:13}\\
Q14 & Ch3, slide 63 & Phosphine cracking: PFR size & p.~\pageref{q:14}\\
Q15 & Ch4, slides 17--20 & Ethylene-oxide batch data: $k$ and batch time & p.~\pageref{q:15}\\
Q16 & Ch4, slides 29--38 & Ethylene glycol: CSTR volume, parallel and series 800-gal reactors & p.~\pageref{q:16}\\
Q17 & Ch4, slides 44--45 & Ethane cracking: PFR volume and number of tubes & p.~\pageref{q:17}\\
Q18 & Ch4, slides 55--58 & Gas $A+B\to2C$ in PBR with pressure drop & p.~\pageref{q:18}\\
Q19 & Ch4, slides 76--77 & PBR + CSTR in series: order, exit pressure, particle size & p.~\pageref{q:19}\\
Q20 & Ch4, slides 78--79 & Series reactions $A\to B\to C$ in a batch reactor & p.~\pageref{q:20}\\
Q21 & Ch4, slides 90--98 & PFR with pressure drop: pipe length for 10\%/20\% & p.~\pageref{q:21}\\
Q22 & Ch5, slides 12--17 & Trityl + methanol: order and rate constant & p.~\pageref{q:22}\\
Q23 & Ch5, slide 75 & Dimethyl ether bomb: order and $k$ & p.~\pageref{q:23}\\
Q24 & Ch5, slide 72 & First-order reversible $A\rightleftharpoons R$ & p.~\pageref{q:24}\\
Q25 & Ch5, slide 73 & Half-order batch; $2A\to B$ volume change & p.~\pageref{q:25}\\
Q26 & Ch5, slide 74 & $A\to3R$ in MFR: rate equation & p.~\pageref{q:26}\\
Q27 & Ch5, slide 76 & $A\to2.5R$ batch, partial-pressure data & p.~\pageref{q:27}\\
Q28 & Ch5, slides 78--79 & Reaction bomb $2A\to B$ at 100\,°C & p.~\pageref{q:28}\\
Q29 & Ch5, slide 80 & Gambling tables (first-order analogue) & p.~\pageref{q:29}\\
Q30 & Ch5, slide 81 & $-r_A=kC_A^2$: 6$\times$ MFR and same-size PFR & p.~\pageref{q:30}\\
Q31 & Ch5, slide 82 & $A+3B\to6R$: $X_A,X_B,C_B$ & p.~\pageref{q:31}\\
Q32 & Ch5, slide 83 & Gaseous feed with $T,\pi$ change: $A+B\to5R$ & p.~\pageref{q:32}\\
Q33 & Ch5, slide 84 & Sucrose--sucrase Michaelis--Menten fit & p.~\pageref{q:33}\\
Q34 & Ch5, slide 61 & Reversible PFR: equilibrium and actual conversion & p.~\pageref{q:34}\\
Q35 & Ch5, slide 62 & $A\to R+2S$ from total-pressure data & p.~\pageref{q:35}\\
Q36 & Ch6, slide 32 & Product distribution, 90\% conversion, three contacting patterns & p.~\pageref{q:36}\\
Q37 & Ch6, slide 33 & Two MFRs in series, $r_R=k_1C_A^2$, $r_S=k_2C_A$ & p.~\pageref{q:37}\\
Q38 & Ch7, slides 14--17 & n-Butane isomerisation: adiabatic PFR volume & p.~\pageref{q:38}\\
Q39 & Ch7, slides 48--51 & Adiabatic equilibrium + two interstage coolers & p.~\pageref{q:39}\\
Q40 & Ch7, slides 57--59 & Propylene-oxide CSTR operated adiabatically & p.~\pageref{q:40}\\
Q41 & Ch7, slides 60--74 & The same CSTR with a cooling coil & p.~\pageref{q:41}\\
Q42 & Ch7, slide 95 & Parallel reactions in a PFR with heat effects & p.~\pageref{q:42}\\
Q43 & Catalysis, slide 48 & Pellet concentration profile; diameter for $\eta=0.8$ & p.~\pageref{q:43}\\
Q44 & Catalysis, slide 57 & Internal diffusion: rates at 100/150\,°C, apparent $E$ & p.~\pageref{q:44}\\
Q45 & Non-ideal, s.17--19 & Tracer pulse: $E(t)$ and fraction between 3 and 6 min & p.~\pageref{q:45}\\
Q46 & Non-ideal, s.52--59 & Dispersion model: conversion four ways & p.~\pageref{q:46}\\
\bottomrule
\end{tabularx}
\end{footnotesize}

\note{The decks are photographs of slides; parts of some questions sit inside the pictures and are quoted as printed.  Decks that repeat a question (the dimerisation of Q13 reappears in Ch5) are solved once and cross-referenced.}

%============================= PART A =====================================
\spsection{Part A -- Mole balances and rates (Ch1-krd (1).pdf)}

\qtitle{Q1 -- Rocket engine: rates of reaction}{\label{q:1}Ch1, slide 7}
\src{Ch1-krd (1).pdf, slide 7 (``Example'').}
A rocket engine burns a stoichiometric mixture of hydrogen and oxygen.  The combustion
chamber is cylindrical, 75\,cm long and 60\,cm in diameter, and produces 108\,kg/s of
exhaust gas.  If combustion is complete, find the rate of reaction of hydrogen and of oxygen.

\begin{qbox}
$\mathrm{H_2 + \tfrac12 O_2 \to H_2O}$.  Chamber volume
$V=\frac{\pi}{4}(0.6)^2(0.75)=0.212\ \mathrm{m^3}$.
Exhaust $=108/18 = 6.0$ kmol/s of $\mathrm{H_2O}$, so 6.0 kmol/s $\mathrm{H_2}$ and
3.0 kmol/s $\mathrm{O_2}$ react (complete combustion).
With $-r_i = F_i/V$:
\[
-r_{\mathrm{H_2}}=\frac{6.0}{0.212}=28.3\ \mathrm{kmol\,m^{-3}\,s^{-1}},\qquad
-r_{\mathrm{O_2}}=\frac{3.0}{0.212}=14.1\ \mathrm{kmol\,m^{-3}\,s^{-1}}.
\]
\end{qbox}
\begin{ansbox}$-r_{\mathrm{H_2}} \approx 2.8\times10^{4}$ and $-r_{\mathrm{O_2}} \approx 1.4\times10^{4}$ mol/(m$^3\cdot$s).\end{ansbox}

\qtitle{Q2 -- CSTR, PFR and batch for three rate laws}{\label{q:2}Ch1, slide 12}
\src{Ch1-krd (1).pdf, slide 12 (``Problem'').}
$A\to B$ isothermally in a continuous-flow reactor; consume 99\% of A
($C_A=0.01\,C_{A0}$) with $F_{A0}=5$ mol/h and $v_0=10$ dm$^3$/h (so $C_{A0}=0.5$ mol/dm$^3$).
Rate laws: (a) $-r_A=k$, $k=0.05$ mol/(h$\cdot$dm$^3$); (b) $-r_A=kC_A$, $k=0.0001$ s$^{-1}=0.36$ h$^{-1}$;
(c) $-r_A=kC_A^2$, $k=3$ dm$^3$/(mol$\cdot$h).  (d) The batch time for 99\% in each case.

\begin{qbox}
CSTR $V=F_{A0}X/(-r_A)_{exit}$; PFR $V=\int_0^X F_{A0}\,dX/(-r_A)$; batch $t=N_{A0}\int_0^X dX/(-r_A V)$ with $C_{A0}=0.5$.
\begin{itemize}[leftmargin=14pt,itemsep=1pt]
\item (a) zero order: $V_{CSTR}=V_{PFR}=5(0.99)/0.05 = 99$ dm$^3$.
\item (b) first order: $V_{CSTR}=\frac{v_0 X}{k(1-X)}=\frac{10(0.99)}{0.36(0.01)}=2750$ dm$^3$;
$V_{PFR}=\frac{v_0}{k}\ln\frac1{1-X}=\frac{10}{0.36}(4.605)=127.9$ dm$^3$.
\item (c) second order: $V_{CSTR}=\frac{v_0 X}{kC_{A0}(1-X)^2}=\frac{10(0.99)}{3(0.5)(10^{-4})}=66\,000$ dm$^3$;
$V_{PFR}=\frac{v_0}{kC_{A0}}\frac{X}{1-X}=\frac{10}{1.5}(99)=660$ dm$^3$.
\item (d) batch: (a) $t=C_{A0}X/k=9.9$ h; (b) $t=\ln(100)/0.36=12.8$ h; (c) $t=\frac{X}{kC_{A0}(1-X)}=66$ h.
\end{itemize}
\end{qbox}
\begin{ansbox}(a) 99 dm$^3$ both; (b) CSTR 2750, PFR 128 dm$^3$; (c) CSTR 66\,000, PFR 660 dm$^3$; (d) 9.9 / 12.8 / 66 h.\end{ansbox}

%============================= PART B =====================================
\spsection{Part B -- Conversion and reactor sizing (Ch2-krd (1).pdf)}

\qtitle{Q3 -- Sizing from the Levenspiel table}{\label{q:3}Ch2, slides 19--24}
\src{Ch2-krd (1).pdf, slides 19--24.}
Calculate the volumes required to achieve 80\% conversion in a CSTR and a PFR when species A
enters at 0.4 mol/s, given $F_{A0}/(-r_A)$ (m$^3$):
\begin{center}\small
\begin{tabular}{lccccccc}
\toprule
$X$ & 0 & 0.1 & 0.2 & 0.4 & 0.6 & 0.7 & 0.8\\
$F_{A0}/(-r_A)$ & 0.89 & 1.08 & 1.33 & 2.05 & 3.54 & 5.06 & 8.00\\
\bottomrule
\end{tabular}
\end{center}
\begin{qbox}
CSTR: $V=\frac{F_{A0}}{-r_A}\Big|_{0.8} X = 8.00(0.8)=6.40$ m$^3$.
PFR with the deck's five-point quadrature ($h=0.2$):
$V=\frac{0.2}{3}[0.89+4(1.33)+2(2.05)+4(3.54)+8.00]=2.165$ m$^3$ (deck value).
\end{qbox}
\begin{ansbox}$V_{CSTR}=6.4$ m$^3$, $V_{PFR}=2.165$ m$^3$.\end{ansbox}

\qtitle{Q4 -- Two CSTRs in series}{\label{q:4}Ch2, slides 31--36}
\src{Ch2-krd (1).pdf, slides 35--36.}
With the same table, the first CSTR reaches $X=0.4$; find the volume of each of two CSTRs
for 80\% overall.
\begin{qbox}
$V_1 = 2.05(0.4)=0.82$ m$^3$; $V_2 = 8.00(0.8-0.4)=3.20$ m$^3$; total $4.02$ m$^3$ ---
less than one CSTR (6.4 m$^3$), as the deck points out.
\end{qbox}
\begin{ansbox}$V_1=0.82$ m$^3$, $V_2=3.20$ m$^3$, total $4.02$ m$^3$.\end{ansbox}

\qtitle{Q5 -- Two PFRs in series}{\label{q:5}Ch2, slides 37--38}
\src{Ch2-krd (1).pdf, slides 37--38.}
Same feed; intermediate conversion 40\%, final 80\%.
\begin{qbox}
$V_1=\frac{0.2}{3}[0.89+4(1.33)+2.05]=0.551$ m$^3$;
$V_2=\frac{0.2}{3}[2.05+4(3.54)+8.00]=1.614$ m$^3$;
total $2.165$ m$^3$ $=$ one PFR, as expected for PFRs in series (deck values).
\end{qbox}
\begin{ansbox}$V_1=0.551$, $V_2=1.614$, total $2.165$ m$^3$.\end{ansbox}

\qtitle{Q6 -- Mixed reactor, unknown stoichiometry}{\label{q:6}Ch2, slides 41--42}
\src{Ch2-krd (1).pdf, slides 41--42.}
A and B ($C_{A0}=0.10$, $C_{B0}=0.01$ mol/L) react in a mixed reactor of $V=1$ L,
$\tau = V/v_0 = 1$ min.  The outlet contains $C_A=0.02$, $C_B=0.03$, $C_C=0.04$ mol/L.
Find the rates of reaction of A, B and C.
\begin{qbox}
$r_i=(C_i-C_{i0})/\tau$: $r_A=-0.08$, $r_B=+0.02$, $r_C=+0.04$ mol/(L$\cdot$min).
Ratio $-r_A:r_B:r_C = 4:1:2$, i.e. the stoichiometry consistent with the data is
$4A \to B + 2C$.
\end{qbox}
\begin{ansbox}$-r_A=0.08$, $r_B=0.02$, $r_C=0.04$ mol/(L$\cdot$min); stoichiometry $4A\to B+2C$.\end{ansbox}

\qtitle{Q7 -- PFRs in two parallel branches}{\label{q:7}Ch2, slides 45--46}
\src{Ch2-krd (1).pdf, slides 45--46.}
Branch D: 20 L then 30 L in series; branch E: 40 L.  What fraction of the feed should go
to branch D?
\begin{qbox}
For parallel PFR branches the efficient split keeps $V/F$ equal:
$F_D/F_E = 50/40$, so $F_D/F = 50/90 = 0.556$.
\end{qbox}
\begin{ansbox}$5/9 \approx 0.56$ of the feed to branch D.\end{ansbox}

\qtitle{Q8 -- P2-7: adiabatic $A+B\to C$, parts (a)--(e)}{\label{q:8}Ch2, slide 59}
\src{Ch2-krd (1).pdf, slide 59 (``P2-7'').}
$F_{A0}=300$ mol/min with
\begin{center}\small
\begin{tabular}{lcccccccc}
\toprule
$X$ & 0 & 0.2 & 0.4 & 0.45 & 0.5 & 0.6 & 0.8 & 0.9\\
$-r_A$ (mol/dm$^3\cdot$min) & 1.0 & 1.67 & 5.0 & 5.0 & 5.0 & 5.0 & 1.25 & 0.91\\
\bottomrule
\end{tabular}
\end{center}
(a) PFR and CSTR volumes for 40\%; (b) range of $X$ where the two volumes are identical;
(c) maximum $X$ in a 105 dm$^3$ CSTR; (d) $X$ for a 72 dm$^3$ PFR followed by a 24 dm$^3$ CSTR;
(e) $X$ for the reverse order.
\begin{qbox}
With $g(X)=F_{A0}/(-r_A) = 300,179.6,60,60,60,60,240,329.7$ dm$^3$:
\begin{itemize}[leftmargin=14pt,itemsep=1pt]
\item (a) CSTR: $60(0.4)=24$ dm$^3$; PFR: trapezium over the table gives $71.9$ dm$^3$.
\item (b) $X\,g(X)=\int_0^X g$ first holds near $X\approx0.69$ (between 0.6 and 0.8, where the
rate falls again); below it the CSTR is smaller, above it the CSTR is larger.
\item (c) $X\,g(X)=105 \Rightarrow X\approx0.70$.
\item (d) PFR to 72 dm$^3$ $\Rightarrow X_1\approx0.40$; then $(X_2-0.4)g(X_2)=24 \Rightarrow X_2\approx0.64$.
\item (e) CSTR: $X_1 g(X_1)=24 \Rightarrow X_1=0.40$ (the high-conversion root); PFR adds
72 dm$^3$ $\Rightarrow X_2\approx0.90$.
\end{itemize}
\end{qbox}
\begin{ansbox}(a) 24 dm$^3$ / 71.9 dm$^3$; (b) about $X=0.69$; (c) $X\approx0.70$; (d) $X\approx0.64$; (e) $X\approx0.90$.\end{ansbox}

%============================= PART C =====================================
\spsection{Part C -- Rate laws and stoichiometry (Ch3-krd.pdf)}

\qtitle{Q9 -- Relative rates: NO oxidation}{\label{q:9}Ch3, slides 5--7}
\src{Ch3-krd.pdf, slides 5--7 (``Example'').}
$2\mathrm{NO}+\mathrm{O_2}\to2\mathrm{NO_2}$; NO$_2$ forms at 4 mol/(m$^3\cdot$s).
\begin{qbox}
$\frac{r_{NO}}{-2}=\frac{r_{O_2}}{-1}=\frac{r_{NO_2}}{2}$ gives
$r_{NO}=-4$ and $r_{O_2}=-2$ mol/(m$^3\cdot$s), exactly as on the deck's solution slide.
\end{qbox}
\begin{ansbox}$r_{\mathrm{NO}}=-4$, $r_{\mathrm{O_2}}=-2$ mol/(m$^3\cdot$s).\end{ansbox}

\qtitle{Q10 -- Relative rates: $2A+3B\to5C$}{\label{q:10}Ch3, slides 6--7}
\src{Ch3-krd.pdf, slides 6--7 (``Exercise'').}
$-r_A = 10$ mol/(dm$^3\cdot$s); find the rates of B and C.
\begin{qbox}
$\frac{r_A}{-2}=\frac{r_B}{-3}=\frac{r_C}{5}$:
$r_B=\frac{-3}{-2}(-10)=-15$; $r_C=\frac{5}{-2}(-10)=+25$ mol/(dm$^3\cdot$s) (deck values).
\end{qbox}
\begin{ansbox}$r_B=-15$, $r_C=+25$ mol/(dm$^3\cdot$s).\end{ansbox}

\qtitle{Q11 -- Saponification stoichiometric table}{\label{q:11}Ch3, slides 42--44}
\src{Ch3-krd.pdf, slides 42--44 (``Example'').}
$3\,\mathrm{NaOH} + (\mathrm{C_{17}H_{35}COO})_3\mathrm{C_3H_5} \to
3\,\mathrm{C_{17}H_{35}COONa} + \mathrm{C_3H_5(OH)_3}$, i.e. $3A+B\to3C+D$, liquid phase,
$X$ = conversion of NaOH.
\begin{qbox}
$\Theta_B=\tfrac13,\ \Theta_C=\Theta_D=0$ (as fed):
$C_A=C_{A0}(1-X)$; $C_B=C_{A0}\left(\Theta_B-\tfrac{X}{3}\right)$;
$C_C=C_{A0}X$; $C_D=C_{A0}\tfrac{X}{3}$ --- the deck's table (slides 43--44), with
$V=V_0$ because the mixture is liquid.
\end{qbox}
\begin{ansbox}Concentrations as above; e.g. with $\Theta_B = N_{B0}/N_{A0}$ the B entry is $C_{A0}(\Theta_B - X/3)$.\end{ansbox}

\qtitle{Q12 -- Flow-system table for $A+2B\to C+D$}{\label{q:12}Ch3, slides 48--51}
\src{Ch3-krd.pdf, slides 48--51 (``Example'').}
Liquid phase, $C_{A0}=1.8$, $C_{B0}=6.6$ kmol/m$^3$; A is the basis.
\begin{qbox}
$\Theta_B=6.6/1.8=3.67$.  Effluent rates: $F_A=F_{A0}(1-X)$, $F_B=F_{A0}(3.67-2X)$,
$F_C=F_{A0}X$, $F_D=F_{A0}X$; concentrations $C_i=F_i/v_0$, e.g.
$C_B=C_{A0}(3.67-2X)$ (deck slides 49--51).
\end{qbox}
\begin{ansbox}$C_A=C_{A0}(1-X)$, $C_B=C_{A0}(3.67-2X)$, $C_C=C_D=C_{A0}X$.\end{ansbox}

\qtitle{Q13 -- Dimerisation $2A\to R$ in a MFR}{\label{q:13}Ch3, slide 61}
\src{Ch3-krd.pdf, slide 61; repeated in Ch5\_-\,krd.pdf, slide 77 (``L 5.2'').}
Pure gaseous A ($C_{A0}=100$ mmol/L) dimerises in a MFR ($V=0.1$ L).  Runs:
$v_0 = 10.0,\,3.0,\,1.2,\,0.5$ L/h with $C_A = 85.7,\,66.7,\,50,\,33.4$ mmol/L.  Find a rate
equation.
\begin{qbox}
Gas phase, $\varepsilon_A=\tfrac{1-2}{2}=-0.5$: $C_A=C_{A0}(1-X)/(1-0.5X)$ gives
$X=0.250,\,0.500,\,0.667,\,0.800$.  MFR design: $-r_A=F_{A0}X/V = v_0 C_{A0}X/V$:
$-r_A = 2502,\,1499,\,800,\,400$ mmol/(L$\cdot$h).
Log--log slope of $(-r_A)$ vs $C_A$ is $1.96\approx2$; $-r_A/C_A^2 = 0.34,\,0.34,\,0.32,\,0.36$.
\end{qbox}
\begin{ansbox}$-r_A \approx 0.34\,C_A^2$ with $C_A$ in mmol/L and $-r_A$ in mmol/(L$\cdot$h) (second order in A).\end{ansbox}

\qtitle{Q14 -- Phosphine cracking PFR}{\label{q:14}Ch3, slide 63}
\src{Ch3-krd.pdf, slide 63 (``Example'').}
$4\,\mathrm{PH_3}\to\mathrm{P_4}+6\,\mathrm{H_2}$ at 649\,°C, first order with
$k=10$ h$^{-1}$.  Size a PFR at 649\,°C and 460 kPa for 80\% conversion of 40 mol/h pure PH$_3$.
\begin{qbox}
$C_{A0}=P/RT=460\,000/(8.314\times922.15)=60.0$ mol/m$^3$; $\varepsilon_A=(10-4)/4=1.5$.
$V=\frac{F_{A0}}{kC_{A0}}\left[(1+\varepsilon)\ln\frac1{1-X}-\varepsilon X\right]
=\frac{40}{10\times60}\left[2.5\ln5-1.2\right]=0.0667(2.419+0.402-0.402)$
$=0.188$ m$^3$.
\end{qbox}
\begin{ansbox}$V\approx0.19$ m$^3$ (about 190 L).\end{ansbox}

%============================= PART D =====================================
\spsection{Part D -- Isothermal design (Ch4-\_Krd.pdf)}

\qtitle{Q15 -- Ethylene-oxide batch data}{\label{q:15}Ch4, slides 17--20}
\src{Ch4-\_Krd.pdf, slides 17--20.}
500 mL of 2 M EO mixed with 500 mL of water (0.9 wt\% H$_2$SO$_4$ catalyst) at 55\,°C
($C_{A0}=1.0$ kmol/m$^3$ after mixing); glycol concentrations recorded vs time.
Determine $k$; and the batch time for 90\% conversion.
\begin{qbox}
$\ln(C_{A0}/C_A)$ vs $t$ is linear; the deck's slope is $k=0.311$ min$^{-1}$ (our least-squares
fit of the printed $\ln$ column gives 0.315).  Then
$t_{0.9}=\ln(10)/k = 7.4$ min.  The deck's scaling slide: with $k=10^{-4}$ s$^{-1}$,
$t_{0.9}=2.3/10^{-4}=23\,000$ s $=6.4$ h.
\end{qbox}
\begin{ansbox}$k\approx0.311$ min$^{-1}$; $t_{90}\approx7.4$ min (and 6.4 h for $k=10^{-4}$ s$^{-1}$).\end{ansbox}

\qtitle{Q16 -- Ethylene glycol CSTR}{\label{q:16}Ch4, slides 29--38}
\src{Ch4-\_Krd.pdf, slides 29--38.}
Produce 200 million lb/yr of EG at 80\% conversion; 1 lbmol/ft$^3$ EO solution mixed with an
equal volume of water (0.9 wt\% H$_2$SO$_4$); $k=0.311$ min$^{-1}$.  (a) CSTR volume;
(b) conversion of two 800-gal reactors in parallel; (c) in series.
\begin{qbox}
$F_C=200\times10^6/62$ lbmol/yr $=3.226\times10^6$/yr; $F_{A0}=F_C/0.8=7.67$ lbmol/min;
$v_{A0}=7.67$ ft$^3$/min, $v_0=15.34$ ft$^3$/min, $C_{A0}=0.5$ lbmol/ft$^3$.
(a) $V=\frac{F_{A0}X}{kC_{A0}(1-X)}=\frac{7.67(0.8)}{0.311(0.5)(0.2)}=197.3$ ft$^3$ (deck value).
(b) Each reactor sees $v_{A0}$: $\tau=\frac{800/7.48}{7.67}=13.94$ min, $Da=k\tau=4.34$,
$X=\frac{Da}{1+Da}=0.81$.
(c) Each sees $v_0$: $\tau=6.97$ min, $Da=2.17$, $X=1-\frac1{(1+Da)^2}=0.90$ (deck values).
\end{qbox}
\begin{ansbox}(a) 197.3 ft$^3$; (b) $X\approx0.81$; (c) $X\approx0.90$.\end{ansbox}

\qtitle{Q17 -- Ethane cracking PFR and tube count}{\label{q:17}Ch4, slides 44--45}
\src{Ch4-\_Krd.pdf, slides 44--45.}
300 million lb/yr ethylene from pure ethane, irreversible elementary rate law, 80\% conversion,
1100 K, 6 atm; $k=0.072$ s$^{-1}$ at 1000 K, $E=82$ kcal/mol.  How many 0.0205 ft$^3$,
40-ft PFR tubes in parallel?
\begin{qbox}
$k(1100)=0.072\exp\!\left[\frac{82\,000}{1.987}\left(\frac1{1000}-\frac1{1100}\right)\right]=3.07$ s$^{-1}$;
$C_{A0}=6\times101\,325/(8.314\times1100)=66.5$ mol/m$^3$;
$F_{A0}=300\times10^6/28$ lbmol/yr $=154$ mol/s; $v_0=2.32$ m$^3$/s.
$\mathrm{C_2H_6\to C_2H_4+H_2}$: $\varepsilon=1$;
$V=\frac{v_0}{k}\left[2\ln\frac1{1-X}-X\right]=\frac{2.32}{3.07}(2.419)=1.83$ m$^3$
$=64.6$ ft$^3$.  Tubes $=64.6/0.0205\approx3.2\times10^3$.
\end{qbox}
\begin{ansbox}$V\approx1.8$ m$^3$ (65 ft$^3$); about 3150 tubes of 0.0205 ft$^3$ in parallel.\end{ansbox}
\begin{warnbox}The deck solves the same problem numerically in MATLAB (slides 84--95) without printing the final volume; the closed-form first-order gas-phase integral above is the method the deck derives on slides 42--43.\end{warnbox}

\qtitle{Q18 -- PBR with pressure drop, $A+B\to2C$}{\label{q:18}Ch4, slides 55--58}
\src{Ch4-\_Krd.pdf, slides 55--58.}
Equimolar feed, $k=1.5$ dm$^6$/(mol$\cdot$kg$\cdot$min), $\alpha=0.0099$ kg$^{-1}$, $\varepsilon=0$.
Find $X$ at $W=100$ kg.
\begin{qbox}
$\frac{dX}{dW}=\frac{kC_{A0}^2}{F_{A0}}(1-X)^2(1-\alpha W)$ integrates to
$\frac{X}{1-X}=\frac{kC_{A0}^2}{F_{A0}}\left(W-\frac{\alpha W^2}{2}\right)$,
$p=(1-\alpha W)^{1/2}$.
The deck prints $X=0.75$ for $\alpha=0$, which fixes the group
$\frac{kC_{A0}^2}{F_{A0}}=\frac{3}{100}=0.03$ kg$^{-1}$; then
$\frac{X}{1-X}=0.03(100-49.5)=1.515 \Rightarrow X=0.60$ --- the deck's two answers.
\end{qbox}
\begin{ansbox}$X=0.60$ with pressure drop, $0.75$ without.\end{ansbox}
\begin{warnbox}The slide does not print $C_{A0}$ and $F_{A0}$ separately; the combined group 0.03 kg$^{-1}$ is implied by its own $\alpha=0$ answer and is stated here instead of being invented.\end{warnbox}

\qtitle{Q19 -- PBR + CSTR in series}{\label{q:19}Ch4, slides 76--77}
\src{Ch4-\_Krd.pdf, slides 76--77.}
First order in $p_A$, $A\to B$ (gas), PBR with 50 kg catalyst gives 50\% conversion at
$P_0=20$ atm; $\alpha=0.018$ kg$^{-1}$.  A CSTR with the same 50 kg is added in series.
(a) Which order gives the higher conversion? (b) The exit conversion; (c) pressure at the
PBR exit; (d) effect of halving catalyst diameter and enlarging the pipe 50\% (turbulent).
\begin{qbox}
Shooting the coupled ODEs $\frac{dX}{dW}=a\frac{1-X}{1+X}p$, $\frac{dp}{dW}=-\frac{\alpha}{2p}$
to match $X(50)=0.5$ gives $a=kP_0/F_{A0}=0.0247$ kg$^{-1}$.
(c) $p=\sqrt{1-\alpha W}=\sqrt{0.1}=0.316 \Rightarrow P=6.3$ atm.
(b) PBR$\to$CSTR: $X=0.60$; CSTR$\to$PBR: $X=0.69$ --- putting the CSTR first is better,
because the PBR then runs at full pressure.
(d) Turbulent $\alpha\propto D_p^{-2}D_{pipe}^{-4}$: $\alpha'=0.018\times2\times1.5^{-4}=0.0071$,
and PBR$\to$CSTR rises to $X\approx0.73$ (similarly higher in the other order).
\end{qbox}
\begin{ansbox}(a) CSTR upstream; (b) $X\approx0.69$ (0.60 the other way); (c) 6.3 atm; (d) $X\approx0.73$.\end{ansbox}

\qtitle{Q20 -- Series reactions in a batch reactor}{\label{q:20}Ch4, slides 78--79}
\src{Ch4-\_Krd.pdf, slides 78--79.}
$A\xrightarrow{k_1}B\xrightarrow{k_2}C$, $k_1=1$ h$^{-1}$, $k_2=2$ h$^{-1}$,
$C_A(0)=5$ mol/m$^3$, $C_B(0)=C_C(0)=0$.
\begin{qbox}
$C_A=5e^{-t}$; $C_B=5\frac{k_1}{k_2-k_1}\left(e^{-k_1t}-e^{-k_2t}\right)=5(e^{-t}-e^{-2t})$;
$C_C=5-C_A-C_B$.  $C_B$ peaks at $t_{max}=\frac{\ln(k_2/k_1)}{k_2-k_1}=0.693$ h with
$C_B^{max}=1.25$ mol/m$^3$.
\end{qbox}
\begin{ansbox}$t_{max}=\ln2\approx0.69$ h, $C_B^{max}=1.25$ mol/m$^3$; profiles as above.\end{ansbox}

\qtitle{Q21 -- PFR pressure drop: pipe length}{\label{q:21}Ch4, slides 90--98}
\src{Ch4-\_Krd.pdf, slides 90--98.}
$A\to B$, first order, $P_0=10$ atm, $Q_0=0.005$ m$^3$/s, 300 K, $k=0.1$ s$^{-1}$,
$M=60$ g/mol, $\mu=10^{-5}$ Pa$\cdot$s.  Length for (a) 2 cm pipe, 10\%; (b) 1.5 cm pipe,
10\%; (c) 1.5 cm pipe, 20\%?
\begin{qbox}
$\rho_0=P_0M/RT=24.4$ kg/m$^3$, $v_0=Q_0/A=15.9$ m/s (2 cm), $F_{A0}=2.03$ mol/s.
Without drop, $L_{10}=\frac{v_0}{k}\ln\frac1{0.9}=16.8$ m and $L_{20}=35.5$ m (2 cm).
The deck integrates the coupled $dX/dz$, $dP/dz$ ODEs (slides 92--95) and reports
(a) $\approx20$ m for 10\% in the 2 cm pipe; (b) 10\% still possible in the 1.5 cm pipe
($\approx$ longer, steeper pressure fall); (c) 20\% \emph{not} possible --- the pressure
collapses before 20\% is reached (slide 98: ``10\% conversion is possible, but 20\% is not'').
\end{qbox}
\begin{ansbox}(a) $\approx20$ m; (b) 10\% achievable; (c) 20\% unattainable in the 1.5 cm pipe.\end{ansbox}

%============================= PART E =====================================
\spsection{Part E -- Collection and analysis of rate data (Ch5\_-\,krd.pdf)}

\qtitle{Q22 -- Trityl + methanol}{\label{q:22}Ch5, slides 12--17}
\src{Ch5\_-\,krd.pdf, slides 12--17 and 42--49.}
$(\mathrm{C_6H_5})_3\mathrm{CCl}+\mathrm{CH_3OH}\to(\mathrm{C_6H_5})_3\mathrm{COCH_3}+\mathrm{HCl}$;
$C_{B0}=0.5$ M; trityl data ($\times10^3$ mol/dm$^3$): 50, 38, 30.6, 25.6, 22.2, 19.5, 17.4 at
$t=0,50,\dots,300$ min.  (1) order in trityl; (2) rate constant given first order in methanol.
\begin{qbox}
Finite-difference rates $-dC_A/dt$ (deck slide 17): $2.86,1.94,1.24,0.84,0.61,0.48,0.36
\times10^{-4}$.  Log--log slope $=1.99\approx2$ (deck) --- second order in trityl.
Integral check: $1/C_A-1/C_{A0}$ vs $t$ is linear with slope $k'=0.124\approx0.125$
dm$^3$/(mol$\cdot$min).  With $-r_A=kC_A^2C_B$ and $C_B\approx C_{B0}$:
$k=k'/0.5=0.25$ dm$^6$/(mol$^2\cdot$min).  The deck's nonlinear regression returns
$\alpha\approx2$ with the same $k'$.
\end{qbox}
\begin{ansbox}Second order in trityl; $k'=0.125$ dm$^3$/(mol$\cdot$min), $k=0.25$ dm$^6$/(mol$^2\cdot$min).\end{ansbox}

\qtitle{Q23 -- Dimethyl ether decomposition bomb}{\label{q:23}Ch5, slide 75}
\src{Ch5\_-\,krd.pdf, slide 75 (``F 5.8'').}
$(\mathrm{CH_3})_2\mathrm{O}\to\mathrm{CH_4}+\mathrm{H_2}+\mathrm{CO}$ at 504\,°C, constant
volume; $\pi$ (mmHg) $=408,488,562,799$ at $t=390,777,1195,3155$ s and $931$ at $t=\infty$.
(a) Why is $\pi_0$ missing, and estimate it; (b) order and rate constant.
\begin{qbox}
(a) The bomb was already hot and reacting before the first reading; since 1 mol $\to$ 3 mol,
$\pi_0=\pi_\infty/3=310$ mmHg.
(b) $p_A=(\pi_\infty-\pi)/2 = 261.5,\,221.5,\,184.5,\,66.0$ mmHg; the four values of
$\ln(p_{A0}/p_A)/t$ are $4.4,\,4.3,\,4.3,\,4.9\times10^{-4}$ s$^{-1}$ --- first order.
\end{qbox}
\begin{ansbox}$\pi_0\approx310$ mmHg; first order, $k\approx4.4\times10^{-4}$ s$^{-1}$.\end{ansbox}

\qtitle{Q24 -- First-order reversible $A\rightleftharpoons R$}{\label{q:24}Ch5, slide 72}
\src{Ch5\_-\,krd.pdf, slide 72.}
$C_{A0}=0.5$ mol/L, $C_{R0}=0$; after 8 min $X=33.3$\%; $X_e=66.7$\%.  Find the rate equation.
\begin{qbox}
$\ln\frac{X_e}{X_e-X}=(k_1+k_2)t$: $\ln 2 = (k_1+k_2)(8)\Rightarrow k_1+k_2=0.0866$ min$^{-1}$;
$K_e=X_e/(1-X_e)=2=k_1/k_2$.  Hence $k_1=0.0578$, $k_2=0.0289$ min$^{-1}$ and
$-r_A=0.0578\,C_A-0.0289\,C_R$ (mol/L$\cdot$min).
\end{qbox}
\begin{ansbox}$-r_A = 0.058\,C_A - 0.029\,C_R$ mol/(L$\cdot$min).\end{ansbox}

\qtitle{Q25 -- Half-order batch; $2A\to B$ volume change}{\label{q:25}Ch5, slide 73}
\src{Ch5\_-\,krd.pdf, slide 73.}
(i) $A\to R$, $-r_A=3C_A^{0.5}$ mol/(L$\cdot$h), $C_{A0}=1$ mol/L: conversion after 1 h?
(ii) Gas $2A\to B$, 80\% A, constant pressure: volume falls 20\% in 3 min; first-order $k$?
\begin{qbox}
(i) $2(\sqrt{C_{A0}}-\sqrt{C_A})=3t$ gives $\sqrt{C_A}<0$ for $t>0.667$ h: the charge is
consumed in 40 min, so $X(1\,\mathrm{h})=1$.
(ii) $V/V_0=1+\varepsilon_A X$ with $\varepsilon_A=0.8\times\frac{1-2}{2}=-0.4$; $0.8=1-0.4X
\Rightarrow X=0.5$; $k=\ln\frac1{1-0.5}/3=0.231$ min$^{-1}$.
\end{qbox}
\begin{ansbox}(i) $X=1$ (complete at 40 min); (ii) $k=0.231$ min$^{-1}$.\end{ansbox}

\qtitle{Q26 -- $A\to3R$ in a MFR (L 5.19)}{\label{q:26}Ch5, slide 74}
\src{Ch5\_-\,krd.pdf, slide 74.}
Pure A at 3 atm, 30\,°C ($C_{A0}=120$ mmol/L) into a 1-L MFR; $v_0=0.06,0.48,1.5,8.1$ L/min
with $C_A=30,60,80,105$ mmol/L.
\begin{qbox}
$\varepsilon_A=2$: $X=\frac{C_{A0}-C_A}{C_{A0}+2C_A}=0.50,0.25,0.143,0.045$;
$-r_A=v_0C_{A0}X/V=3.6,14.4,25.7,44.0$ mmol/(L$\cdot$min).  Log--log slope $=2.00$;
$-r_A/C_A^2=0.0040$ in all four runs.
\end{qbox}
\begin{ansbox}$-r_A=0.004\,C_A^2$ ($C_A$ in mmol/L, $-r_A$ in mmol/(L$\cdot$min)), i.e. $4$ L/(mol$\cdot$min) in molar units.\end{ansbox}

\qtitle{Q27 -- $A\to2.5R$ batch, pressure data (L 3.25)}{\label{q:27}Ch5, slide 76}
\src{Ch5\_-\,krd.pdf, slide 76.}
0\,°C, constant volume, pure A: $p_A$ (mm) $=760,600,475,380,320,275,240,215,150$ at
$t=0,2,\dots,14,\infty$.
\begin{qbox}
$p_{A,\infty}=150\neq0$: the reaction approaches equilibrium.  First order towards it:
$\ln\frac{p_{A0}-p_{Ae}}{p_A-p_{Ae}}=kt$ gives $k=0.152$--$0.162$ at every point
(fit $0.160$ min$^{-1}$) --- a first-order reversible rate equation represents the data.
\end{qbox}
\begin{ansbox}First-order approach to $p_{Ae}=150$ mm with $k\approx0.16$ min$^{-1}$.\end{ansbox}

\qtitle{Q28 -- Reaction bomb $2A\to B$ at 100\,°C}{\label{q:28}Ch5, slides 78--79}
\src{Ch5\_-\,krd.pdf, slides 78--79.}
Pure A at 1 atm and 25\,°C is jumped to 100\,°C; $\pi$ falls from 1.14 atm at $t=1$ min
to 0.728 atm at $t=20$ min (twelve readings, e.g. 1.04, 0.982, 0.940, 0.905, 0.870 at
$t=2..6$); after the weekend no A remains.  Rate equation in mol, L, min.
\begin{qbox}
$2A\to B$ with $\pi_0(100\,°\mathrm{C})=1$ atm: $p_A=2\pi-1$.  Least squares:
second order ($1/p_A$ linear) fits with $r^2=0.964$ vs $0.886$ first order, giving
$k_p=0.074$ atm$^{-1}$min$^{-1}$; with $p_A=C_ART$ ($RT=30.6$ L$\cdot$atm/mol at 373 K),
$k_C=k_pRT\approx2.3$ L/(mol$\cdot$min).
\end{qbox}
\begin{ansbox}Second order in A; $k\approx0.074$ atm$^{-1}$min$^{-1}$ $\approx2.3$ L/(mol$\cdot$min) at 100\,°C.\end{ansbox}

\qtitle{Q29 -- The gambling tables}{\label{q:29}Ch5, slide 80}
\src{Ch5\_-\,krd.pdf, slide 80.}
Losses proportional to money held: table 1 halves it in 4 h, table 2 in 2 h.  Starting with
Rs.\,1000 at both tables at once, when is Rs.\,10 left?
\begin{qbox}
First-order decay with $k=\ln2/4+\ln2/2=0.520$ h$^{-1}$;
$t=\ln(1000/10)/k=4.605/0.520=8.9$ h.
\end{qbox}
\begin{ansbox}About 8.9 hours.\end{ansbox}

\qtitle{Q30 -- Second order: bigger MFR, equal PFR}{\label{q:30}Ch5, slide 81}
\src{Ch5\_-\,krd.pdf, slide 81.}
$-r_A=kC_A^2$ gives 50\% in a MFR.  (i) conversion in a MFR of 6 times the volume;
(ii) in a PFR of the same size.
\begin{qbox}
Original: $kC_{A0}\tau=\frac{X}{(1-X)^2}=2$.
(i) $kC_{A0}\tau'=12=\frac{X}{(1-X)^2}\Rightarrow X=0.75$.
(ii) PFR: $kC_{A0}\tau=\frac{X}{1-X}=2\Rightarrow X=2/3$.
\end{qbox}
\begin{ansbox}(i) $X=0.75$; (ii) $X=0.67$.\end{ansbox}

\qtitle{Q31 -- $A+3B\to6R$ with inerts}{\label{q:31}Ch5, slide 82}
\src{Ch5\_-\,krd.pdf, slide 82.}
Feed $C_{A0}=100$, $C_{B0}=200$, $C_{I0}=100$; $C_A=40$ at the exit.  Find $X_A,X_B,C_B$.
\begin{qbox}
$X_A=0.6$; $X_B=X_A\frac{3C_{A0}}{C_{B0}}\cdot\frac{1}{3}\times3=0.9$; $C_B=200-3(60)=20$ mol/L.
\end{qbox}
\begin{ansbox}$X_A=0.6$, $X_B=0.9$, $C_B=20$.\end{ansbox}

\qtitle{Q32 -- Gaseous feed with $T$ and $\pi$ change}{\label{q:32}Ch5, slide 83}
\src{Ch5\_-\,krd.pdf, slide 83.}
$T_0=1000$ K, $\pi_0=5$ atm, $C_{A0}=100$, $C_{B0}=200$; $A+B\to5R$; at the exit
$T=400$ K, $\pi=4$ atm, $C_A=20$.  Find $X_A,X_B,C_B$.
\begin{qbox}
$y_{A0}=1/3$, $\delta=+3$, $\varepsilon_A=1$:
$C_A=C_{A0}\frac{1-X}{1+X}\frac{T_0}{T}\frac{\pi}{\pi_0}=32\frac{1-X}{1+X}=20
\Rightarrow X_A=0.231$; $X_B=X_A C_{A0}/C_{B0}=0.115$;
$C_B=200\frac{1-X_B}{1+X_A}\times0.32=46.0$.
\end{qbox}
\begin{ansbox}$X_A=0.23$, $X_B=0.12$, $C_B\approx46$.\end{ansbox}

\qtitle{Q33 -- Sucrase kinetics, Michaelis--Menten}{\label{q:33}Ch5, slide 84}
\src{Ch5\_-\,krd.pdf, slide 84.}
$C_{A0}=1.0$, $C_{S0}=0.01$ mmol/L; $C_A$ (mmol/L) $=0.68,0.60,0.55,0.38,0.27,0.16,0.09,0.04,
0.018,0.006,0.0025$ at $t=1..11$ h.  Does Michaelis--Menten fit?  Evaluate $k_3,C_M$.
\begin{qbox}
Integrated MM: $k_3C_{S0}\,t=(C_{A0}-C_A)+C_M\ln\frac{C_{A0}}{C_A}$, linearised as
$\frac{t}{\ln(C_{A0}/C_A)}=\frac{1}{k_3C_{S0}}\frac{C_{A0}-C_A}{\ln(C_{A0}/C_A)}+\frac{C_M}{k_3C_{S0}}$.
Least squares over the eleven points: slope $=3.12$ h, intercept $=1.61$ h, giving
$k_3C_{S0}=0.321$ h$^{-1}$ ($k_3\approx32$ h$^{-1}$) and $C_M\approx0.52$ mmol/L; the points
lie on the line ($r^2>0.99$), so the fit is reasonable.
\end{qbox}
\begin{ansbox}Yes: $k_3\approx32$ h$^{-1}$ (at $C_{S0}=0.01$), $C_M\approx0.5$ mmol/L.\end{ansbox}

\qtitle{Q34 -- Reversible PFR: equilibrium then actual conversion}{\label{q:34}Ch5, slide 61}
\src{Ch5\_-\,krd.pdf, slide 61.}
A PFR (2 m$^3$) processes aqueous feed (100 L/min) with $C_{A0}=100$ mmol/L; $A\rightleftharpoons B$.
Find the equilibrium conversion, then the actual conversion (and repeat for a CSTR).
\begin{qbox}
The slide prints no rate constants; the textbook problem it reproduces (Levenspiel, the
companion MFR version of which is printed with $-r_A=0.04C_A-0.01C_R$ mol/(L$\cdot$min))
supplies them: $K_e=k_1/k_2=4\Rightarrow X_e=\frac{4}{5}=0.80$.
$-r_A=0.05(C_A-0.02)$; $\tau=V/v_0=20$ min;
$\ln\frac{C_{Af}-0.02}{0.08}=-(k_1+k_2)\tau=-1\Rightarrow C_{Af}=0.0494$, $X=0.51$.
(CSTR: $0.05(C_A-0.02)\times20=C_{A0}-C_A\Rightarrow C_A=0.06$, $X=0.40$.)
\end{qbox}
\begin{ansbox}$X_e=0.80$; PFR $X\approx0.51$; CSTR $X\approx0.40$ (rate law quoted from the textbook original, as the slide omits it).\end{ansbox}

\qtitle{Q35 -- $A\to R+2S$ from total-pressure data}{\label{q:35}Ch5, slide 62}
\src{Ch5\_-\,krd.pdf, slide 62.}
Pure A, constant volume; $\pi$ (mmHg) $=7.5,10.5,12.5,15.8,17.9,19.4$ at
$t=0,2.5,5,10,15,20$ min.
\begin{qbox}
$\Delta n=+2$: $p_A=\pi_0-\frac{\pi-\pi_0}{2}=\frac{3\pi_0-\pi}{2}=7.5,6.0,5.0,3.35,2.3,1.55$.
$\ln(p_{A0}/p_A)/t = 0.089,0.081,0.081,0.079,0.079$ --- first order, $k\approx0.08$ min$^{-1}$.
\end{qbox}
\begin{ansbox}First order with $k\approx0.08$ min$^{-1}$.\end{ansbox}

%============================= PART F =====================================
\spsection{Part F -- Multiple reactions (Ch6\_krd.pdf)}

\qtitle{Q36 -- Product distribution, three contacting patterns}{\label{q:36}Ch6, slide 32}
\src{Ch6\_krd.pdf, slide 32.}
Aqueous $A+B\xrightarrow{k_1}R$ (desired), $A+B\xrightarrow{k_2}S$, with
$r_R=1.0\,C_A^{1.5}C_B^{0.3}$, $r_S=1.0\,C_A^{0.5}C_B^{1.8}$ mol/(L$\cdot$min).
Equal streams of A and B at 20 mol/L each are fed ($C_{A0}=C_{B0}=10$ after mixing);
90\% conversion of A.  Find $C_R$ for (a) plug flow, (b) mixed flow, (c) plug flow with side
entrances of B.
\begin{qbox}
With $C_A=C_B=c$: $\varphi(R/A)=\frac{r_R}{r_R+r_S}=\frac{1}{1+c^{0.5}}$.
(a) $C_R=\int_{1}^{10}\frac{dc}{1+\sqrt c}=2\left[u-\ln(1+u)\right]_1^{\sqrt{10}}=2.86$ mol/L.
(b) $C_R=\varphi|_{c=1}\,(10-1)=0.5\times9=4.5$ mol/L.
(c) Side entry of B holds $C_B$ at the low exit level ($=1$):
$\varphi=\frac{c}{c+1}$, $C_R=\int_1^{10}\frac{c\,dc}{c+1}=9-\ln\frac{11}{2}=7.30$ mol/L.
\end{qbox}
\begin{ansbox}(a) 2.86; (b) 4.50; (c) 7.30 mol/L --- mixed flow beats plug flow here, and keeping B low is best.\end{ansbox}

\qtitle{Q37 -- Two MFRs in series}{\label{q:37}Ch6, slide 33}
\src{Ch6\_krd.pdf, slide 33.}
$r_R=k_1C_A^2$, $r_S=k_2C_A$; feed $C_{A0}=1.0,C_{R0}=0,C_{S0}=0.3$; $\tau_1=2.5$, $\tau_2=10$ min;
first-reactor exit $C_A=0.4,C_R=0.2,C_S=0.7$.  Find the second-reactor exit.
\begin{qbox}
From reactor 1: $k_1=\frac{0.2/2.5}{0.4^2}=0.5$; $k_2=\frac{(0.7-0.3)/2.5}{0.4}=0.4$
(the S balance checks: $0.4\times0.4=0.16=(0.7-0.3)/2.5$).
Reactor 2: $0.4-C_A=(0.5C_A^2+0.4C_A)10\Rightarrow C_A=0.0745$;
$C_R=0.2+0.5C_A^2(10)=0.228$; $C_S=0.7+0.4C_A(10)=0.998$.
\end{qbox}
\begin{ansbox}$C_A\approx0.075$, $C_R\approx0.23$, $C_S\approx1.00$ mol/L.\end{ansbox}

%============================= PART G =====================================
\spsection{Part G -- Non-isothermal design (Ch7-\_krd.pdf)}

\qtitle{Q38 -- Adiabatic PFR: n-butane isomerisation}{\label{q:38}Ch7, slides 14--17}
\src{Ch7-\_krd.pdf, slides 14--17.}
$n$-butane $\rightleftharpoons$ isobutane, adiabatic liquid-phase PFR; $k=31.1$ h$^{-1}$ at
360 K, $E=65.7$ kJ/mol, $\Delta H_{Rx}=-6900$ J/mol, $K_C=3.03$ at 60\,°C,
$C_{A0}=9.3$ kmol/m$^3$, $C_p=141/141/161$ J/(mol$\cdot$K) (butane/isobutane/pentane).
Process 163 kmol/h (90 mol\% butane) to $X=0.70$, feed at 330 K.
\begin{qbox}
Adiabatic line $T=330+43.3X$ (deck); $k(T)=31.1e^{7902(1/360-1/T)}$;
$K_C(T)=3.03e^{830.3(1/333-1/T)}$; $-r_A=kC_{A0}\left[1-\left(1+\tfrac1{K_C}\right)X\right]$;
$F_{A0}=0.9\times163=146.7$ kmol/h.  The deck's table (e.g. $-r_A=39.2$ at $X=0$,
$F_{A0}/-r_A=23.29$ m$^3$ at 0.7) is reproduced by these expressions; integrating
$F_{A0}/(-r_A)$ to $X=0.7$ gives $V\approx2.0$ m$^3$ (the deck's coarse Simpson table reads
$\approx2.4$ m$^3$; the value is sensitive near $X_e\approx0.72$).
\end{qbox}
\begin{ansbox}$V\approx2$--$2.4$ m$^3$ for 70\% (deck table: $F_{A0}/(-r_A)=23.3$ m$^3$ at $X=0.7$).\end{ansbox}

\qtitle{Q39 -- Adiabatic equilibrium with two interstage coolers}{\label{q:39}Ch7, slides 48--51}
\src{Ch7-\_krd.pdf, slides 48--51.}
$A\rightleftharpoons B$, pure A at 300 K; $K_C(298)=10^5$, $\Delta H_{Rx}=-20\,000$ cal/mol
(as implied by the deck's $K_C(T)$ and adiabatic line), $C_{pA}=C_{pB}=50$ cal/(mol$\cdot$K),
so $X_{ad}=2.5\times10^{-3}(T-300)$.  Two coolers return the stream to 350 K; 95\% of the
equilibrium conversion in each reactor; $F_{A0}=40$ mol/s; cooling water at 270 K,
$c_{pC}=18$ cal/(mol$\cdot$K), $U=100$ cal/(s$\cdot$m$^2\cdot$K).
\begin{qbox}
$K_C(T)=10^5\exp[-33.78(T-298)/T]$, $X_e=K_C/(1+K_C)$.
Reactor 1: $2.5\times10^{-3}(T-300)=0.95X_e(T)\Rightarrow T_1=459$ K, $X_1=0.40$.
Cool to 350 K.  Reactor 2: $X_1+2.5\times10^{-3}(T-350)=0.95X_e(T)\Rightarrow T_2=440$ K,
$X_2=0.62$.
Duties: $\dot Q=\dot n C_p (T_{out}-350)$: $Q_1=40\times50\times109=218$ kcal/s;
$Q_2=40\times50\times90=179$ kcal/s (water flow $=\dot Q/(18\,\Delta T)$ as sized by the deck).
\end{qbox}
\begin{ansbox}$X\approx0.40$ after reactor 1 and $0.62$ after reactor 2; $\dot Q_1\approx218$, $\dot Q_2\approx179$ kcal/s.\end{ansbox}

\qtitle{Q40 -- Propylene-oxide CSTR, adiabatic}{\label{q:40}Ch7, slides 57--59}
\src{Ch7-\_krd.pdf, slides 57--59 (and 69--72).}
300-gal glass-lined CSTR; $F_{A0}=43.04$ lbmol/h PO; streams: 46.62 ft$^3$/h PO, 46.62 ft$^3$/h
methanol, 233.1 ft$^3$/h water (0.1 wt\% H$_2$SO$_4$); methanol 71.87 and water 802.8 lbmol/h;
$T_0=75$\,°F $=535$ R after the 17\,°F heat-of-mixing rise; first order in PO with
$k=16.96\times10^{12}e^{-32\,400/RT}$ h$^{-1}$; limit 125\,°F.
\begin{qbox}
$v_0=326.3$ ft$^3$/h, $C_{A0}=0.132$ lbmol/ft$^3$, $\tau=(300/7.48)/326.3=0.123$ h.
Energy balance: $\sum\Theta_iC_{p,i}\approx404$ Btu/(lbmol A$\cdot$R), $-\Delta H_{Rx}=36\,400$
Btu/lbmol $\Rightarrow T=535+90.1X$.  Solving with $X=\frac{k\tau}{1+k\tau}$:
$X=0.84$, $T=611$ R $=151$\,°F (deck table intersection $\approx613$ R, $X\approx0.85$).
$151>125$\,°F: the reactor \emph{cannot} be used adiabatically.
\end{qbox}
\begin{ansbox}$X\approx0.84$--$0.85$ at $T\approx611$--$613$ R ($\approx151$\,°F) $>125$\,°F --- not acceptable.\end{ansbox}

\qtitle{Q41 -- The same CSTR with a cooling coil}{\label{q:41}Ch7, slides 60--74}
\src{Ch7-\_krd.pdf, slides 60--74.}
A 40 ft$^2$ coil, $U=100$ Btu/(h$\cdot$ft$^2\cdot$°F), coolant at 85\,°F (545 R).  Does the
125\,°F constraint now hold?
\begin{qbox}
Energy balance adds $UA(T-T_a)=4000(T-545)$ Btu/h:
$X=\frac{17\,379(T-535)+4000(T-545)}{43.04\times36\,400}$ together with the mole balance.
Intersection: $X\approx0.35$--$0.36$, $T\approx562$--$564$ R $=102$--$104$\,°F --- the deck's
fsolve gives $X=0.364$, $T=563.7$ R.  The constraint is satisfied with the coil.
\end{qbox}
\begin{ansbox}Yes: $X\approx0.36$ at $T\approx564$ R ($\approx104$\,°F) $<125$\,°F.\end{ansbox}

\qtitle{Q42 -- Parallel reactions in a PFR with heat effects}{\label{q:42}Ch7, slide 95}
\src{Ch7-\_krd.pdf, slide 95.}
$A\to B$ ($r_B=k_1C_A$), $2A\to C$ ($r_C=k_2C_A^2$); pure A at 100 mol/s, 150\,°C,
$C_{A0}=0.1$ mol/dm$^3$; $\Delta H_1=-20\,000$, $\Delta H_2=-60\,000$ J/mol A;
$C_{pA}=C_{pB}=90$, $C_{pC}=180$ J/(mol$\cdot$°C); $E_1/R=4000$ K, $E_2/R=9000$ K;
$k_1=10$, $k_2=0.09$ (as printed, per min).  Determine the temperature and flow-rate
profiles down the reactor.
\begin{qbox}
Coupled ODEs (per the deck's PFR heat-effects balances):
$\frac{dF_A}{dV}=-(r_1+2r_2)$, $\frac{dF_C}{dV}=r_2$,
$\frac{dT}{dV}=\frac{r_1(-\Delta H_1)+r_2(-\Delta H_2)}{F_AC_{pA}+F_BC_{pB}+F_CC_{pC}}$
with $C_A=F_A/\dot V_0$, $\dot V_0=60\,000$ dm$^3$/min (taking the printed rate constants as
per minute).  Integrating to $V=10$ m$^3$: $F_A\to0$, $F_C\to0$ (B dominates), and $T$ rises
to $\approx645$ K --- the adiabatic limit for complete conversion to B
($\Delta T = 100\times20\,000/(100\times90)\approx222$ K above 423 K).
\end{qbox}
\begin{ansbox}Both reactions complete within the first m$^3$; $T\to\approx645$ K, product essentially B.\end{ansbox}
\begin{warnbox}The slide states the problem and the data but no units for $k_1,k_2$ beyond ``min''; the profile above takes them as printed (per minute).  The qualitative result (fast exothermic run to the adiabatic temperature rise) is unaffected.\end{warnbox}

%============================= PART H =====================================
\spsection{Part H -- Catalysis (Catalysis\_Krd.pdf)}

\qtitle{Q43 -- Pellet concentration profile}{\label{q:43}Catalysis, slide 48}
\src{Catalysis\_Krd.pdf, slide 48 (``Problem'').}
$A\to B$, first order, spherical particle $d_p=2\times10^{-3}$ cm; $C_{As}=0.001$ gmol/m$^3$;
$D_e=0.1$ cm$^2$/s.  The concentration halfway to the centre is $1/10$ of $C_{As}$.
(a) $C_A$ at $3\times10^{-4}$ cm from the external surface; (b) diameter for $\eta=0.8$.
\begin{qbox}
Sphere, first order: $\frac{C_A}{C_{As}}=\frac{R}{r}\frac{\sinh(\phi r/R)}{\sinh\phi}$.
At $r=R/2$: $0.1=\frac{2\sinh(\phi/2)}{\sinh\phi}=\frac1{\cosh(\phi/2)}\Rightarrow
\phi=2\,\mathrm{arcosh}(10)=5.99$.
(a) $r=R-3\times10^{-4}=7\times10^{-4}$ cm: ratio $=\frac{1}{0.7}\frac{\sinh(0.7\times5.99)}{\sinh 5.99}=0.237$,
$C_A=2.4\times10^{-4}$ gmol/m$^3$.
(b) $\eta=\frac{3}{\phi^2}(\phi\coth\phi-1)=0.8\Rightarrow\phi_{0.8}=2.04$; $\phi\propto d_p$:
$d_{new}=2\times10^{-3}\times\frac{2.04}{5.99}=6.8\times10^{-4}$ cm.
\end{qbox}
\begin{ansbox}(a) $C_A\approx2.4\times10^{-4}$ gmol/m$^3$; (b) $d_p\approx6.8\times10^{-4}$ cm.\end{ansbox}

\qtitle{Q44 -- Internal diffusion: 100 vs 150\,°C}{\label{q:44}Catalysis, slide 57}
\src{Catalysis\_Krd.pdf, slide 57.}
First order, $R=0.2$ cm, $D_e=0.015$ cm$^2$/s; $k(100\,°\mathrm{C})=0.93$ s$^{-1}$,
$E=20$ kcal/mol; $C_{Ab}=3.25\times10^{-5}$ mol/L.  (a) rate at 100\,°C; (b) at 150\,°C;
(c) apparent activation energy.
\begin{qbox}
(a) $\phi_1=R\sqrt{k_1/D_e}=1.575$; $\eta_1=0.866$; $-r_A=\eta_1 k_1 C_{Ab}=2.6\times10^{-5}$
mol/(L$\cdot$s).
(b) $k_2=0.93e^{\frac{20\,000}{1.987}(1/373-1/423)}=22.5$ s$^{-1}$; $\phi_2=7.75$;
$\eta_2=0.337$; $-r_A=2.5\times10^{-4}$ mol/(L$\cdot$s).
(c) From the two rates, $E_{app}=R\ln(r_2/r_1)/(1/T_1-1/T_2)\approx14.1$ kcal/mol --- between
$E$ and $E/2$, moving toward the diffusion-limited $E/2=10$ kcal/mol as $\phi$ grows.
\end{qbox}
\begin{ansbox}(a) $2.6\times10^{-5}$; (b) $2.5\times10^{-4}$ mol/(L$\cdot$s); (c) $E_{app}\approx14$ kcal/mol.\end{ansbox}

%============================= PART I =====================================
\spsection{Part I -- Residence time distributions (Non\_ideal\_KRD.pdf)}

\qtitle{Q45 -- Tracer pulse: $E(t)$ and the 3--6 min fraction}{\label{q:45}Non-ideal, slides 17--19}
\src{Non\_ideal\_KRD.pdf, slides 17--19 (``Problem'').}
$C$ (g/m$^3$) $=0,1,5,8,10,8,6,4,3,2.2,1.5,0.6,0$ at $t=0,1,\dots,10,12,14$ min.  Construct
$E(t)$ and find the fraction that spent 3--6 min in the reactor.
\begin{qbox}
$\int_0^\infty C\,dt = 47.4+2.6 = 50.0$ g$\cdot$min/m$^3$ (Simpson on 0--10 plus 10--14, as on
the deck).  $E(t)=C(t)/50$: peak $E(4)=0.20$.  Fraction between 3 and 6 min:
$\int_3^6 E\,dt \approx \frac{1}{50}\times\frac{3}{8}(8+3\cdot10+3\cdot8+6)=0.465\approx0.5$.
Mean residence time $t_m=\int tE\,dt\approx5.2$ min.
\end{qbox}
\begin{ansbox}$\int C\,dt=50$ g$\cdot$min/m$^3$; about half the material (0.47--0.51) spent 3--6 min inside; $t_m\approx5.2$ min.\end{ansbox}

\qtitle{Q46 -- Dispersion model: conversion four ways}{\label{q:46}Non-ideal, slides 52--59}
\src{Non\_ideal\_KRD.pdf, slides 52--59 (``Problem'').}
First order $A\to B$ in a 10 cm, 6.36 m tubular reactor, $k=0.25$ min$^{-1}$; tracer test of
Q45.  Conversion by (i) closed-vessel dispersion model, (ii) PFR, (iii) single CSTR,
(iv) tanks in series.
\begin{qbox}
$t_m=5.13$ min, $\sigma^2=5.92$ min$^2$ (deck values), $\sigma_\theta^2=0.225$;
$\sigma_\theta^2=\frac{2}{Pe}-\frac{2}{Pe^2}(1-e^{-Pe})\Rightarrow Pe=7.7$.
$Da=k t_m=1.28$, $q=\sqrt{1+4Da/Pe}=1.29$:
$X=1-\frac{4q\,e^{Pe/2}}{(1+q)^2e^{Pe\,q/2}-(1-q)^2e^{-Pe\,q/2}}=0.68$ (deck value).
(ii) $X=1-e^{-Da}=0.72$; (iii) $X=\frac{Da}{1+Da}=0.56$;
(iv) $n=t_m^2/\sigma^2=4.45$: $X=1-(1+Da/n)^{-n}=0.68$.
\end{qbox}
\begin{ansbox}(i) 0.68; (ii) 0.72; (iii) 0.56; (iv) 0.68.\end{ansbox}

\vspace{6pt}
\noindent\rule{\linewidth}{0.5pt}\\[3pt]
{\small Every number above is recomputed by \texttt{krd-answers-check.py} (kept in this
folder) from the data quoted here; discrepancies with a deck's own printed answer, where any,
are discussed in that question's Notes/Warning box.}

\end{document}
